\documentclass{report}

\input{../templates/preamble}
\input{../templates/macros}
\input{../templates/letterfonts}

\title{\Huge{}\ undertext}
\author{\Huge{Marcus Allen Denslow}}
\date{}

\begin{document}

\maketitle
\newpage% or \cleardoublepage
% \pdfbookmark[<level>]{<title>}{<dest>}
\pdfbookmark[section]{\contentsname}{toc}
\tableofcontents
\pagebreak

\chapter{Chapter Title}
\section{Section Title}

\chapter{Alle burde bytte til Linux}
  Me har alle vore der, me skal til å gjere noko ekstremt viktig på pcen, men så frys han.
  Stressfølelsen vert overvelda av sinne, alt arbeidet er kanskje borte, dei 4 timane du
  brukte på å skrive ein stil på skriveøkta i nynorsk ligg i hendene til Windows. Du held
  inne av knappen for å starte PC-en på nytt, du ventar litt med ei forventning om at
  skjermen blir svart, men det blir ho ikkje. Du stussar litt, tenkjer "kva i all verda er
  det som skjer". Det er nemleg ei oppdatering på PC-en. Oppdateringa er berre 10 minutt,
  men det er nok til å øydeleggje alt, den svellande kjensla av håplausheit over kor
  kontrollaus du er tek over.

  Windows har lenge vorte brukt som det primære operativsystemet i skuleverket. Du får
  faktisk ikkje lov til å bruke PC utan at han køyrer Windows i norske skular. Skule-pcen
  har Windows, og vidare ein custom Windows-versjon for å hindre din eigen kontroll over
  din PC. Som ein som bruker skule-pc meiner eg at me har rett til å velje eit
  operativsystem som passar oss. Eg vil bringe opp ei hending som skjedde med meg under ein
   prøve i Norsk i VG2. Eg sat nervøst og venta på å få lov til å opne pcen, deretter opne
  Word og skrive noko som William Shakespeare ville vore stolt av. Men tida frå å trykkje
  på på-knappen til å opne Word var over 40 minutt. 

  Dei 40 minutta var halve prøvetida. Når operativsystemet hindrar elevar frå å prestere på
   skulen, har eleven rett på ny PC eller fridom til å velje eit operativsystem som gjer
  PC-en brukbar.

  Tilbake til eksempelet eg brukte i starten, så har du heller ikkje noko kontroll over kva
   Windows gjer og når. Windows bryr seg ikkje om du ikkje gjer noko eller skriv eksamen.
  Når Windows vil oppdatere, så oppdaterer han, ingen spørsmål spurte. Dette gjer ikkje
  Linux. Der har ein full kontroll over når oppdateringar skjer, så om det verkeleg ikkje
  passar no, så ikkje oppdater no. Det er så enkelt.

notater:
hendene, bruker, eksempel, 
ikkje bruk så mange kommaer.



\end{documentment}
